\begin{figure*}
  \begin{minipage}{0.49\textwidth}
  \[
  \begin{array}{lrcl}
    \text{stacks}
        & \Stk{K} & \Coloneqq\  & \Stk{\StkEmpty}         \\
        &         &             & \Stk{\StkKont{x}{N}{K}} \\
        &         &             & \Stk{\StkVal{V}{K}}     \\[0.8ex]
    \text{logic ctx.}
        & \Gamma  & \Coloneqq\  & \\
        &         &             & \Gamma, x : \sigma      \\[0.8ex]
    \text{susp.\ env.}
        & \Susp{S} & \Coloneqq & \SuspEmpty{} \\
        &          &           & \AddSusp{S}{x}{M}
  \end{array}
  \]
\end{minipage}
\begin{minipage}{0.49\textwidth}
  \[
  \begin{array}{lrcl}
    \text{basic machine}
        & \Machine{L} & \Coloneqq & \SLMachine{M}{K} \\[0.8ex]
    \text{configurations}
        & \Conf{C} & \Coloneqq & \FailC{}       \\
        &          &           & \Machine{L}    \\
        &          &           & \ChoiceC{C}{C} \\[0.8ex]
    \text{terminals}
        & \CTerm{T} & \Coloneqq\ & \CTerm{\Return{V}}     \\
        &           &            & \CTerm{\Lam{x}{M}}
  \end{array}
  \]
\end{minipage}

\bigskip
\raggedright

\fbox{ $\MSteps{L}{C}\phantom{'} $ } \quad \textbf{Basic transitions} \\[1ex]

\emph{Structural transitions}:

\begin{minipage}{0.48\textwidth}
\begin{align*}
  \SLMachine{\CTo{M}{x}{N}}{K}
    &\MStep{} \SLMachine{N}[\AddSusp{S}{x}{M}]{K}
\end{align*}
\end{minipage}
\begin{minipage}{0.48\textwidth}
\begin{align*}
  \SLMachine{\Return{V}}{\StkKont{x}{N}{K}}
    &\MStep{} \SLMachine{\CSubst{N}{V}{x}}{K}
\end{align*}
\end{minipage}

\begin{align*}
  \SLMachine{M}[\DisjSusp{S}{x}{N}]{K}
    &\MStep \SLMachine{N}{\StkKont{x}{M}{K}}
    && \text{(if $x$ is needed in $\CTerm{M}$)} \\[0.5ex]
  \SLMachine{\CApp{M}{V}}{K}
    &\MStep{} \SLMachine{M}{\StkVal{V}{K}} \\
  \SLMachine{\Lam{x}{M}}{\StkVal{V}{K}}
    &\MStep{} \SLMachine{\CSubst{M}{V}{x}}{K} \\
  \SLMachine{\Force*{\Thunk{M}}}{K}
    &\MStep \SLMachine{M}{K} \\
  \SLMachine{\Fix{x}{M}}{K}
    &\MStep \SLMachine{\CSubst{M}{\Thunk*{\Fix{x}{M}}}{x}}{K} \\
  \SLMachine{\CIfz{\Zero}{M}{n}{N}}{K}
    &\MStep \SLMachine{M}{K} \\
  \SLMachine{\CIfz{\Succ{V}}{M}{n}{N}}{K}
    &\MStep \SLMachine{\CSubst{N}{V}{n}}{K} \\
  \SLMachine{\CMatch{\Nil}{M}{x}{xs}{N}}{K}
    &\MStep \SLMachine{M}{K} \\
  \SLMachine{\CMatch{\Cons{V}{W}}{M}{x}{xs}{N}}{K}
    &\MStep \SLMachine{\SubstTwo{N}{V}{W}{x}{xs}}{K}
\end{align*}

\emph{Effect transitions}: \hfill
\begin{minipage}{0.9\textwidth}
\begin{align*}
  \SLMachine{\Fail}{K}
    &\MStep{} \FailC{} \\
  \SLMachine{\Choice{M}{N}}{K}
    &\MStep{} \ChoiceC{\SLMachine{M}{K}}{\SLMachine{N}{K}} \\
  \SLMachine{\Exists{x}[\sigma]{M}}{K}
    &\MStep{} \SLMachine[\Gamma, x : \sigma]{M}{K}
\end{align*}
\end{minipage}

\emph{Narrowing transitions}
  (only if $\SLMachine{M}{K}$ is \emph{blocked} on $x$):
\begin{align*}
  \SLMachine[\Gamma, x : \NatT, \Gamma']{M}{K}
    &\MStep
    \ChoiceC{
      \MSubst{\SLMachine[\Gamma, \Gamma']{M}{K}}{\Zero}{x}
    }{
      \MSubst{\SLMachine[\Gamma, n : \NatT, \Gamma']{M}{K}}{\Succ{n}}{x}
    } \\
  \SLMachine[\Gamma, x : \ListT{\sigma}, \Gamma']{M}{K}
    &\MStep
    \ChoiceC{
      \MSubst{\SLMachine[\Gamma, \Gamma']{M}{K}}{\Nil}{x}
    }{
      \MSubst{\SLMachine[\Gamma, y : \sigma, ys :
      \ListT{\sigma}, \Gamma']{M}{K}}{\Cons{y}{ys}}{x}
    }
\end{align*}

\emph{Unification transitions}:
\begin{align*}
  \SLMachine[\Gamma]{\Equate{V}{W}{M}}{K}
    &\MStep \SimSubM{\prn{\SLMachine[\Gamma]{M}{K}}}{\gamma}
    && (\Mgu{\Gamma,\, \VTerm{V},\, \VTerm{W}} = (\Delta,\, \gamma)) \\
  \SLMachine[\Gamma] {\Equate{V}{W}{M}}{K}
    &\MStep \FailC{}
    && (\Mgu{\Gamma,\, \VTerm{V},\, \VTerm{W}} = \bot)
\end{align*}

\fbox{ $\CSteps{C}{C'}$ } \quad \textbf{Configuration transitions}
\begin{mathpar}
  \inferH{Leaf}{
    \MSteps{L}{C} \\
  }{
    \CSteps{\Ctx{M}[L]}{\Ctx{M}[C]}
  }
  \and
  \inferH{Refl}{
    \strut
  }{
    \CSteps{C}{C}
  }
  \and
  \inferH{Trans}{
    \CSteps{C_1}{C_2} \\
    \CSteps{C_2}{C_3}
  }{
    \CSteps{C_1}{C_3}
  }
\end{mathpar}

  \caption{Abstract Machine. A basic machine $\SLMachine{M}{K}$ carries a
    logic context $\Gamma$, a computation $\CTerm{M}$, a suspension environment
    $\Susp{S}$, and a stack $\Stk{K}$. The bind transition suspends
    $\CTerm{M}$ in $\Susp{S}$ rather than evaluating it immediately; the
    forcing transition reverses this when $x$ is needed. Narrowing transitions
    fire only when a logical variable blocks the current computation.}
  \label{figure:machine}
\end{figure*}

The CBPV term produced by the translation of \cref{section:ir} is not evaluated
directly. Instead, it gives rise to a \emph{nondeterministic search tree}: a
(possibly infinite) tree whose internal nodes are choice points and whose leaves
are either results or failures. The task of execution is to explore this tree
and collect the results. We decompose this task into two parts: a
\emph{basic machine} that steps a single CBPV computation until it reaches a
branch point, and a \emph{search strategy} that manages the collection of
active machines.

This decomposition is a direct realisation of the abstract machine of
\parencite{jones_2026}. That paper organises machine states into
\emph{configurations}: binary trees whose leaves are basic machines or
failures. A configuration transition selects a non-terminal leaf and steps it,
possibly splitting it into further branches. The theory is silent on
\emph{which} leaf to step:
%
\begin{quote}
  \emph{``Observe that this is highly nondeterministic: it may be that one or
  more leaves of $\Conf{C}$ may be able to make a basic transition! We leave
  the choice to the implementer, who may use any strategy to explore the tree
  (e.g.~DFS, BFS, parallel).''} \hfill---\parencite{jones_2026}
\end{quote}
%
\noindent This freedom is justified by the denotational semantics: because the
semantics is based on the lower powerdomain, which is insensitive to the order
of exploration, all strategies produce the same set of results for terminating
programs.

\subsection{The Basic Machine}

The basic machine (\cref{figure:machine}) is a state machine whose transitions
correspond to the evaluation of a single CBPV computation. Its state consists
of:
%
\begin{itemize}
  \item A \emph{computation closure} $(\CTerm{M}, \rho)$: the current
    computation paired with an environment $\rho$ mapping de Bruijn indices to
    value closures.
  \item A \emph{stack} $\Stk{K}$: a list of frames representing pending
    continuations ($\textsf{to}\ x.\, \CTerm{N}$), argument values
    ($\VTerm{V}$), and suspension-set frames that record which suspension to
    update when a forced computation completes.
  \item A \emph{logic context} $\Gamma$: a record of logical variables, their
    first-order types, and their current bindings.
  \item A \emph{suspension environment} $\Susp{S}$: a map from suspension
    identifiers to deferred computation closures (discussed in
    \cref{section:suspensions}).
\end{itemize}

\noindent The machine steps by pattern-matching on the current computation. Most
transitions are standard CBPV reductions: a $\CTerm{\Return{V}}$ pops a
continuation from the stack and substitutes $\VTerm{V}$ into the continuation
body; an application $\CTerm{\CApp{M}{V}}$ pushes $\VTerm{V}$ onto the stack;
forcing a thunk $\CTerm{\Force*{\Thunk{M}}}$ unwraps the suspended computation;
a case expression selects the appropriate branch based on the head constructor
of the scrutinee. Each of these transitions corresponds to a $\beta$-law of the
equational theory (\cref{figure:equations}).

The machine is structured around a tight inner loop that executes deterministic
steps until a branch point (choice or narrowing) or a terminal state (return
with empty stack, or failure) is reached. Control is returned to the search
strategy only at points where a genuine scheduling decision must be made; all
deterministic transitions are executed in place.

\subsection{Effect Transitions}

The FLP effects produce the branching structure of the search tree.

\paragraph{Choice}
When the machine encounters $\CTerm{\Choice{M}{N}}$, it produces a
\emph{branch point}: two copies of the current state, one continuing with
$\CTerm{M}$ and the other with $\CTerm{N}$. This is the fundamental source of
nondeterminism. The equational theory guarantees that choice is associative,
commutative, and idempotent, so the tree's branching structure does not affect
the set of results.

\paragraph{Failure}
$\CTerm{\Fail}$ discards the current branch: it is a leaf of the search tree
with no result.

\paragraph{Logical variables}
$\CTerm{\Exists{x}[\sigma]{M}}$ allocates a fresh entry in the logic context
with type $\sigma$ and no binding, extends the environment with a reference to
this entry, and continues with $\CTerm{M}$. The variable $x$ is initially
unconstrained; its value is determined later by narrowing or unification.

\paragraph{Narrowing}
The critical interaction between logical variables and case analysis is handled
by narrowing. When the machine encounters a case expression whose scrutinee
resolves---through the environment and logic context---to an \emph{unbound}
logical variable $x : \sigma$, it cannot proceed: the variable's head
constructor is unknown. The machine responds by \emph{splitting} into branches,
one for each constructor of $\sigma$. In each branch, $x$ is bound to the
appropriate constructor applied to fresh logical variables.

For example, if $x : \NatT$ is unbound and the machine reaches
$\textsf{ifz}(x; \ldots)$, it produces two branches: one binding $x$ to
$\VTerm{\Zero}$, and one binding $x$ to $\VTerm{\Succ{y}}$ for a fresh
$y : \NatT$. This corresponds to the \emph{parameter equation}
$\CTerm{\Exists{n}[\NatT]{M}} =
\CTerm{\Choice{\CSubst{M}{\Zero}{n}}{\Exists{m}[\NatT]{\CSubst{M}{\Succ{m}}{n}}}}$
of \cref{figure:equations}: the equation says that introducing a logical
variable of type $\NatT$ is the same as choosing between zero and the successor
of a fresh variable.

The machine narrows only when a variable \emph{blocks} a computation, in the
sense of \parencite{jones_2026}: the next step directly depends on the
variable's value. This strategy, called \emph{weak head narrowing}, ensures that
narrowing happens only when necessary. Its completeness---every closed value of
type $\sigma$ is eventually considered as a possible instantiation---is proved
in \parencite{jones_2026}.

\paragraph{Unification}
$\CTerm{\Equate{V}{W}{M}}$ attempts to unify $\VTerm{V}$ and $\VTerm{W}$ using
the standard first-order unification algorithm. The algorithm processes a queue
of pairs of value closures, decomposing constructors and binding variables via
the logic context. Three outcomes are possible: \emph{success}, in which the
logic context is updated with the most general unifier and execution continues
with $\CTerm{M}$; \emph{failure}, in which the values have incompatible head
constructors and the branch is discarded; or \emph{suspension}, in which one of
the values refers to a deferred computation that must be forced before
unification can proceed.

An optional occurs check prevents the construction of infinite terms (e.g.\
$x = \textsf{S}(x)$). It can be disabled for performance at the cost of
soundness.

\subsection{Search Strategies}
\label{subsection:search-strategies}

The search tree generated by the machine must be explored to collect results.
\Gweek{} implements four strategies, each corresponding to a different
traversal of the tree.

\paragraph{Breadth-first search (\BFS{})}
The default strategy maintains two vectors: the current generation and the next
generation. Each machine in the current generation is stepped to its next branch
point; the resulting branches are placed in the next generation. When the
current generation is exhausted, the vectors are swapped. \BFS{} is
\emph{complete}---if a solution exists at any finite depth, it will be
found---and \emph{fair}: no branch is starved. Its limitation is memory: the
number of active machines can grow exponentially in the depth of the tree.

\paragraph{Depth-first search (\DFS{})}
\DFS{} maintains a stack and always steps the most recently created branch. It
uses memory proportional to the depth of the tree, but it is
\emph{incomplete}: if the tree has an infinite branch, \DFS{} may follow it
forever. Infinite branches arise naturally when logical variables of type
$\NatT$ or $\ListT{A}$ are narrowed, so \DFS{} is appropriate only for finite
search spaces or when a single solution is sufficient.

\paragraph{Iterative deepening (\IDDFS{})}
Iterative deepening combines the completeness of \BFS{} with the memory
efficiency of \DFS{}. The strategy runs \DFS{} with a depth limit that starts
at~1 and doubles after each complete pass. Solutions found in previous passes
are deduplicated by a hash set of previously seen outputs. The overhead of
re-exploring shallow nodes is bounded by a constant factor.

\paragraph{Fair round-robin (\Fair{})}
The fair strategy maintains a queue of \emph{threads}, where each thread is a
local \DFS{} work stack. A single thread is created initially; new threads are
created at branch points. When the machine encounters a branch while executing a
thread, the first alternative continues in the current thread and each remaining
alternative is placed in the queue as a new single-machine thread. Each thread
is allowed a fixed quota of steps (currently $10{,}000$). If a thread exhausts
its quota without completing, its remaining work stack is re-enqueued as a
single entry, and the scheduler moves to the next thread.

The key invariant is that \emph{branch alternatives are distributed across
threads}, so every branch receives its own share of the round-robin schedule.
Within a thread, execution is depth-first, so the local work stack remains
small (proportional to the depth of the sub-tree being explored). To bound
memory, the total number of threads is capped at a configurable maximum
(currently $10{,}000$); when the cap is reached, new branch alternatives remain
in their parent thread and are explored depth-first. This degrades gracefully:
the branches that have already been distributed across threads continue to
receive fair scheduling, while only new branches within a full thread lose
individual time-slicing.

This strategy is both complete and fair: every branch is eventually explored,
and no branch is starved indefinitely. It avoids the memory blow-up of \BFS{}
while retaining completeness, and is particularly effective for programs that
mix finite branching (explicit choice) with infinite branching (logical variable
narrowing).

\paragraph{Discussion}
The theoretical licence to choose any exploration order is powerful, but it
places a practical burden on the programmer. In general, \BFS{} is the safest
default for programs involving logical variables, as it does not require the
programmer to reason about the finiteness of branches. \DFS{} is fastest for
finite search spaces. The fair strategy offers a practical compromise when the
structure of the search space is unknown.
